\chapter{Implementation}
\label{ch:implementation}

This chapter describes the implementation part of the project. Before starting to create a project in Firebase, make sure its package name matches the project package name created in Android Studio. This chapter will provide major logic algorithms to help build a complete app.

\section{Algorithms}

% Algorithms 1-21

\textbf{Algorithm 1: Splash Screen}\\
\textit{Input:} App starts, user login status\\
\textit{Output:} User is sent to Home screen or Login screen\\
\textit{Steps:}
\begin{enumerate}
    \item App opens and splash screen is shown.
    \item App checks if the user is already logged in.
    \item If user is logged in, open Home screen.
    \item If user is not logged in, open Login screen.
    \item Splash screen closes.
\end{enumerate}

\textbf{Algorithm 2: Sign Up}\\
\textit{Input:} Name, email, date of birth, gender, password\\
\textit{Output:} New account is created, user goes to Login screen\\
\textit{Steps:}
\begin{enumerate}
    \item User opens Sign Up screen.
    \item User enters name, email, DOB, gender, and password. When user clicks Doctor, Caregiver, Pharmacist option, they enter license number or speciality and send request to admin for approval/rejection.
    \item App checks all fields are filled correctly.
    \item If something is wrong, show error message.
    \item If all correct, create the account.
    \item Send a verification email to the user.
    \item Take user to Login screen.
\end{enumerate}

\textbf{Algorithm 3: Login}\\
\textit{Input:} Email and password\\
\textit{Output:} User is logged in and goes to Home screen\\
\textit{Steps:}
\begin{enumerate}
    \item User opens Login screen.
    \item User enters email and password.
    \item App checks that fields are not empty.
    \item If empty, show error message.
    \item App checks email and password are correct.
    \item If email is not verified, ask user to verify first.
    \item If correct, take user to Home screen.
    \item If user forgot password, send a reset link to their email.
\end{enumerate}

\textbf{Algorithm 4: Patient Module}\\
\textit{Input:} Patient login, health task details\\
\textit{Output:} Patient manages all health tasks in one place\\
\textit{Steps:}
\begin{enumerate}
    \item Patient logs in and opens Home screen.
    \item Patient can add medication reminders.
    \item Patient can set reminders for BP, sugar, and weight.
    \item Patient can set reminders for daily activities.
    \item Patient can book an appointment with a doctor.
    \item Patient can view doctor recommendations.
    \item Patient can mark tasks as taken or missed.
    \item Patient can view health history in Profile.
\end{enumerate}

\textbf{Algorithm 5: Doctor Module}\\
\textit{Input:} Doctor login, patient profile\\
\textit{Output:} Doctor views patient data and gives recommendations\\
\textit{Steps:}
\begin{enumerate}
    \item Doctor logs in with email and password.
    \item Doctor opens the patient list.
    \item Doctor selects a patient to view.
    \item Doctor sees patient's medicines, measurements, and reminder history.
    \item Doctor checks if patient is taking medicines on time.
    \item Doctor writes recommendations like medicine changes or checkup schedule.
    \item Doctor decides if patient needs a caregiver or can manage alone.
    \item Patient can see doctor's notes in their profile.
    \item Doctor logs out when done.
\end{enumerate}

\textbf{Algorithm 6: In-App Chat}\\
\textit{Input:} Approved request status, text message\\
\textit{Output:} Message exchange and payment detail sharing\\
\textit{Steps:}
\begin{enumerate}
    \item Doctor and caregiver approves patient request.
    \item App updates status to Approved.
    \item App enables Chat button for users.
    \item User opens chat and types a message.
    \item If any field is missing, display error message.
    \item If valid, send and display message on screens.
\end{enumerate}

\textbf{Algorithm 7: Payment}\\
\textit{Input:} Payment details (Doctor name, payment number, Transaction ID, receipt screenshot)\\
\textit{Output:} Payment marked as Approved or Rejected\\
\textit{Steps:}
\begin{enumerate}
    \item Doctor opens the Payments tab in the Doctor Portal.
    \item App shows all pending patient payments.
    \item Doctor selects a payment to view the receipt screenshot and Transaction ID.
    \item Doctor taps Approve or Reject.
    \item App updates the payment status in the database.
    \item Patient sees the updated status on their screen.
\end{enumerate}

\textbf{Algorithm 8: Caregiver Module}\\
\textit{Input:} Caregiver login, patient profile\\
\textit{Output:} Caregiver monitors and manages patient health tasks\\
\textit{Steps:}
\begin{enumerate}
    \item Caregiver logs in with email and password.
    \item Caregiver opens assigned patient's profile.
    \item Caregiver can see all active reminders of the patient.
    \item Caregiver updates task status as taken or missed on behalf of patient.
    \item Caregiver tracks patient's daily health progress.
    \item Caregiver views patient's measurement records.
    \item Caregiver makes sure patient does not miss any health task.
    \item Caregiver logs out when done.
\end{enumerate}

\textbf{Algorithm 9: Medication Reminder}\\
\textit{Input:} Medicine name, dosage, frequency, start date, end date, reminder time\\
\textit{Output:} Medication reminder is saved and notification is set\\
\textit{Steps:}
\begin{enumerate}
    \item User opens Medication screen.
    \item User enters medicine name, dosage, frequency, and reminder time.
    \item App checks all fields are filled.
    \item If not filled, show error.
    \item If filled, save the medication reminder.
    \item Set a notification at the reminder time.
    \item Show medicine in the reminder list.
    \item User marks medicine as taken or missed when notification comes.
    \item User can update or delete the reminder anytime.
\end{enumerate}

\textbf{Algorithm 10: Measurement Reminder}\\
\textit{Input:} Measurement type (BP / Sugar / Weight), reminder time\\
\textit{Output:} Measurement reminder is saved and notification is set\\
\textit{Steps:}
\begin{enumerate}
    \item User opens Measurement screen.
    \item User selects measurement type and sets reminder time.
    \item App checks that type and time are selected.
    \item If not selected, show error.
    \item If selected, save the reminder.
    \item Set a notification at the reminder time.
    \item Show reminder in active list.
    \item User can update or delete the reminder anytime.
\end{enumerate}

\textbf{Algorithm 11: Activity Reminder}\\
\textit{Input:} Activity type (Water / Meal / Walk / Exercise), reminder time, frequency\\
\textit{Output:} Activity reminder is saved and notifications are set\\
\textit{Steps:}
\begin{enumerate}
    \item User opens Activity screen.
    \item User selects activity type, reminder time, and frequency.
    \item App checks all options are selected.
\end{enumerate}

\textbf{Algorithm 12: Medical Professional}\\
\textit{Input:} Connection request, Appointment schedule\\
\textit{Output:} Patient connects with doctor/caregiver and appointments saved with reminder\\
\textit{Steps:}
\begin{enumerate}
    \item User opens the Medical Professional screen.
    \item User enters doctor details, reasons/symptoms and sends request.
    \item User enters caregiver name and email and sends request.
    \item App checks that all fields are filled.
    \item If any field is empty, shows error message.
    \item App sends the request and updates status to Pending.
    \item User selects Appointment tab to view or manage scheduled visits.
\end{enumerate}

\textbf{Algorithm 13: Appointment Reminder}\\
\textit{Input:} Appointment date, time, doctor or clinic name\\
\textit{Output:} Appointment is saved and reminder notification is set\\
\textit{Steps:}
\begin{enumerate}
    \item User opens Appointment screen.
    \item User enters date, time, and doctor or clinic name.
    \item App checks all fields are filled and date is in the future.
    \item If not, show error.
    \item If correct, save the appointment.
    \item Set a reminder notification at the appointment time.
    \item Show appointment in upcoming list.
    \item User can update or delete the appointment anytime.
\end{enumerate}

\textbf{Algorithm 14: Refill Reminder}\\
\textit{Input:} Medicine name, quantity\\
\textit{Output:} Refill reminder is saved, notification sent when medicine is low\\
\textit{Steps:}
\begin{enumerate}
    \item User opens medication screen.
    \item User enters medicine name and current quantity.
    \item App checks fields are filled.
    \item If not, show error.
    \item When quantity becomes low, send a notification to refill.
    \item Show refill reminder in active list.
    \item User can update or delete the refill reminder.
\end{enumerate}

\textbf{Algorithm 15: Measurement Tracker}\\
\textit{Input:} Measurement type, values (BP readings / sugar level / weight)\\
\textit{Output:} Health value is saved and shown in history\\
\textit{Steps:}
\begin{enumerate}
    \item User opens Measurement Tracker screen.
    \item User selects measurement type.
    \item User enters the value (e.g., systolic and diastolic for BP).
    \item App checks values are valid numbers.
    \item If not valid, show error.
    \item If valid, save the record with date and time.
    \item Show record in health history list.
    \item User can update or delete any record.
\end{enumerate}

\textbf{Algorithm 16: Pharmacy}\\
\textit{Input:} Medicine name, quantity, delivery address\\
\textit{Output:} Checked medicine stock, calculated total price, placed delivery order\\
\textit{Steps:}
\begin{enumerate}
    \item User opens Pharmacy screen.
    \item User enters medicine name and required quantity.
    \item User taps check delivery availability.
    \item App checks stock and delivery availability.
    \item If out of stock, app displays an error message.
    \item If available, patient enters delivery address.
    \item Pharmacist calculates medicine cost and delivery fee.
    \item Patient confirms and places order using Cash on Delivery.
    \item App saves order details and displays order confirmation to patient.
\end{enumerate}

\textbf{Algorithm 17: Chatbot}\\
\textit{Input:} User's health question\\
\textit{Output:} AI gives a simple health answer or tip\\
\textit{Steps:}
\begin{enumerate}
    \item User opens Chatbot screen.
    \item User types a health question.
    \item Question is sent to the AI model.
    \item AI reads the question and finds the best answer.
    \item If user shares a health value like high BP, AI gives a helpful tip.
    \item AI reply is shown in the chat.
    \item If AI does not understand, it shows a helpful default message.
    \item Chat history stays on screen during the session.
\end{enumerate}

\textbf{Algorithm 18: Notification and Reminder}\\
\textit{Input:} Reminder time set by user in any module\\
\textit{Output:} Notification is sent, status marked as Taken or Missed\\
\textit{Steps:}
\begin{enumerate}
    \item App checks all saved reminders.
    \item At the set time, a notification is sent to the user.
    \item Notification shows the task and reminder type.
    \item If user marks it as Taken, status is saved as Taken.
    \item If user ignores it, notification is repeated until user marks it.
    \item Result is saved in Track Intake.
\end{enumerate}

\textbf{Algorithm 19: Track Intake}\\
\textit{Input:} User response to reminder (Taken or Missed)\\
\textit{Output:} Daily health record is updated\\
\textit{Steps:}
\begin{enumerate}
    \item User opens Track Intake screen.
    \item All today's reminders are shown.
    \item User selects a reminder and marks it as Taken or Missed.
    \item Status is saved with date and time.
    \item User can see daily summary of all Taken and Missed tasks.
    \item All records are saved and can be viewed in Profile History.
\end{enumerate}

\textbf{Algorithm 20: Profile Section}\\
\textit{Input:} User profile info, feedback text, emergency contact details\\
\textit{Output:} Profile updated, feedback sent to admin, emergency contact saved\\
\textit{Steps:}
\begin{enumerate}
    \item User opens Profile screen.
    \item User can see their personal details like name and email.
    \item User goes to History to see all Taken and Missed records.
    \item User goes to Feedback and writes a suggestion or complaint.
    \item App checks feedback is not empty.
    \item If not empty, feedback is saved and sent to admin.
    \item User adds emergency contact.
    \item App checks name and number are valid.
    \item Emergency contact is saved for quick access in emergencies.
\end{enumerate}

\textbf{Algorithm 21: Admin Module}\\
\textit{Input:} Admin login credentials\\
\textit{Output:} Admin can view registered users and user feedback\\
\textit{Steps:}
\begin{enumerate}
    \item Admin enters email and password to login.
    \item Admin opens Admin Dashboard.
    \item Admin can see list of all registered users.
    \item Admin can approve and reject license of doctor.
    \item Admin can view all feedback sent by users.
    \item Admin can view emergency contact details saved by users.
    \item Admin logs out when done.
\end{enumerate}


\section{External APIs/SDKs}

Medicare uses Firebase services for its main external services. Firebase Authentication is used for user registration while Cloud Firestore is used for data management and Firebase Cloud Messaging for push notifications. Flutter Local Notification is used to display scheduled or instant notifications on a user's device. All these services are included via Gradle dependencies and are used as built-in components within the Android project.

\section{Code Repository}

Git repository link: \url{https://github.com/laibashamim189-dotcom/medicaree}. The source code of the project is stored in a Git repository, which contains all the modules and features of Medicare. It is a public repository; the entire code is there and will be updated on a daily basis as we test our app and improve it.

\section{Summary}

This chapter outlined almost all major component algorithms, although not all are included here. One can get the understanding, working, and flow of the things from the algorithms and understand the project from the repository.
